\documentclass[aspectratio=169]{beamer}
\usepackage[utf8]{inputenc}
\usepackage[T1]{fontenc}
\usepackage{graphicx}
\usepackage{caption}
\usepackage{booktabs}
\usepackage{hyperref}
\usepackage{lmodern}
\usepackage{tikz}
\usetikzlibrary{backgrounds,calc}

\usetheme{Madrid}
\useoutertheme{infolines}
\usecolortheme{seahorse}
\setbeamertemplate{navigation symbols}{}
\setbeamertemplate{footline}[frame number]
\setbeamertemplate{itemize items}[circle]
\setbeamersize{text margin left=8mm,text margin right=8mm}
\setbeamercolor{frametitle}{bg=}
\setbeamertemplate{frametitle}{\vspace{2pt}\begin{flushleft}\insertframetitle\end{flushleft}\vspace{-4pt}}
\setbeamercolor{normal text}{fg=white,bg=}
\setbeamercolor{frametitle}{fg=white,bg=}

\definecolor{ffBase}{RGB}{20,24,26}
\definecolor{ffGreenSoft}{RGB}{64,90,78}
\definecolor{ffGreenStrong}{RGB}{47,158,104}
\definecolor{ffGraphite}{RGB}{214,223,218}
\definecolor{ffBlueGlow}{RGB}{98,130,210}

\newcommand{\backgroundcanvas}{%
  \begin{tikzpicture}[remember picture,overlay]
    \path [fill=ffBase,opacity=0.96]
      (current page.south west) rectangle (current page.north east);
    \fill [ffGreenSoft,opacity=0.35] ($(current page.north east)+(-1.1,-0.9)$) circle (3.0);
    \fill [ffBlueGlow!32,opacity=0.30] ($(current page.south west)+(0.8,0.6)$) circle (2.6);
  \end{tikzpicture}%
}
\setbeamertemplate{background}{\backgroundcanvas}

\graphicspath{{dark/}}
\title{Flowforge UI \textendash{} system low-code do budowy interfejsów API}
\subtitle{Tryb ciemny}
\author{Konrad Firlej \\ Promotor: dr inż. Krzysztof Augustynek}
\date{Luty 2026}
\titlegraphic{\vspace{-6mm}\includegraphics[height=1.4cm]{Uniwersytet_Bielsko-Bialski_-_logo.jpg}}

\newcommand{\shot}[2]{%
  \begin{center}
    \includegraphics[height=0.7\textheight,keepaspectratio]{#1}\\[4pt]
    \parbox{0.9\linewidth}{\centering\small #2}
  \end{center}
}

\begin{document}

\begin{frame}
  \titlepage
\end{frame}

\begin{frame}{Agenda}
  \begin{itemize}
    \item Widoki główne (Workflows)
    \item Akcje na workflow (More \textrightarrow{} Versions/Rename/Import)
    \item Konfiguracja bloków (If, Switch, Calculation, HTTP itp.)
    \item Edytor workflow (canvas, palette, variables)
    \item Run drawer (open, result)
    \item Widoki główne (Executions)
    \item Detale wykonania (Execution details)
    \item Widoki główne (Scheduler)
    \item Harmonogram (create/edit)
    \item Widoki główne (Blocks)
  \end{itemize}
\end{frame}

% --- Widoki główne (Workflows) ---
\begin{frame}{Workflows}
  \shot{workflows-dark.png}{Lista projektów workflow \textendash{} tryb ciemny}
\end{frame}

% --- Akcje workflow ---
\begin{frame}{Workflow actions \textendash{} More}
  \shot{workflows-actions-01b-more-open.png}{Menu akcji \textquotedbl More\textquotedbl{} (ciemny)}
\end{frame}

\begin{frame}{Workflow actions \textendash{} Versions}
  \shot{workflows-actions-02-versions.png}{Lista wersji workflow (ciemny)}
\end{frame}

\begin{frame}{Workflow actions \textendash{} Rename}
  \shot{workflows-actions-03-rename-form.png}{Formularz zmiany nazwy (ciemny)}
\end{frame}

\begin{frame}{Workflow actions \textendash{} Import}
  \shot{workflows-actions-04-import.png}{Import workflow z pliku (ciemny)}
\end{frame}

% --- Konfiguracja bloków ---
\begin{frame}{If block}
  \shot{workflow-editor-config-if.png}{Konfiguracja bloku If \textendash{} tryb ciemny}
\end{frame}

\begin{frame}{Switch block}
  \shot{workflow-editor-config-switch.png}{Konfiguracja bloku Switch \textendash{} tryb ciemny}
\end{frame}

\begin{frame}{Calculation block}
  \shot{workflow-editor-config-calculation.png}{Konfiguracja bloku Calculation \textendash{} tryb ciemny}
\end{frame}

\begin{frame}{Text Transform block}
  \shot{workflow-editor-config-texttransform.png}{Konfiguracja bloku Text Transform \textendash{} tryb ciemny}
\end{frame}

\begin{frame}{Text Replace block}
  \shot{workflow-editor-config-textreplace.png}{Konfiguracja bloku Text Replace \textendash{} tryb ciemny}
\end{frame}

\begin{frame}{HTTP request block}
  \shot{workflow-editor-config-httprequest.png}{Konfiguracja bloku HTTP request \textendash{} tryb ciemny}
\end{frame}

\begin{frame}{Parser block}
  \shot{workflow-editor-config-parser.png}{Konfiguracja bloku Parser \textendash{} tryb ciemny}
\end{frame}

\begin{frame}{Loop block}
  \shot{workflow-editor-config-loop.png}{Konfiguracja bloku Loop \textendash{} tryb ciemny}
\end{frame}

\begin{frame}{Wait block}
  \shot{workflow-editor-config-wait.png}{Konfiguracja bloku Wait \textendash{} tryb ciemny}
\end{frame}

% --- Edytor workflow ---
\begin{frame}{Workflow editor}
  \shot{workflow-editor.png}{Początkowy widok edytora \textendash{} tryb ciemny}
\end{frame}

\begin{frame}{Palette}
  \shot{workflow-editor-add-blocks.png}{Paleta dodawania bloków \textendash{} tryb ciemny}
\end{frame}

\begin{frame}{Variables}
  \shot{workflow-editor-variables.png}{Lista zmiennych (ciemny)}
\end{frame}

\begin{frame}{Variables \textendash{} edycja}
  \shot{workflow-editor-variables-edit.png}{Edycja zmiennej (ciemny)}
\end{frame}

% --- Run drawer ---
\begin{frame}{Run \textendash{} panel danych}
  \shot{run-drawer-open.png}{Panel Run z rozwiniętymi danymi wejściowymi i snippetem API}
\end{frame}

\begin{frame}{Run \textendash{} wynik uruchomienia}
  \shot{run-drawer-result.png}{Panel Run po wykonaniu workflow \textendash{} widoczne wyniki i logi}
\end{frame}

% --- Widoki główne (Executions) ---
\begin{frame}{Executions}
  \shot{executions-dark.png}{Historia wykonania skryptów \textendash{} tryb ciemny}
\end{frame}

% --- Execution details ---
\begin{frame}{Execution details}
  \shot{execution-details.png}{Pełne szczegóły przebiegu \textendash{} tryb ciemny}
\end{frame}

% --- Widoki główne (Scheduler) ---
\begin{frame}{Scheduler}
  \shot{scheduler-dark.png}{Lista zaplanowanych zadań \textendash{} tryb ciemny}
\end{frame}

% --- Scheduler actions ---
\begin{frame}{Scheduler action \textendash{} create}
  \shot{scheduler-actions-01-create.png}{Tworzenie wpisu schedulera \textendash{} tryb ciemny}
\end{frame}

\begin{frame}{Scheduler action \textendash{} edit}
  \shot{scheduler-actions-02-edit.png}{Edycja istniejącego wpisu \textendash{} tryb ciemny}
\end{frame}

% --- Widoki główne (Blocks) ---
\begin{frame}{Blocks}
  \shot{blocks-dark.png}{Widok bloków systemowych \textendash{} tryb ciemny}
\end{frame}

\end{document}
